% 本文件由 paper/generate-content.js 从代表性提示词自动生成，请勿手工编辑。
\begingroup
\setlength{\tabcolsep}{0pt}
\renewcommand{\arraystretch}{1.08}
\setlength{\emergencystretch}{3em}
\begin{longtable}{@{}>{\RaggedLeft\arraybackslash\footnotesize}p{0.04\linewidth}@{\hspace{0.8em}}>{\RaggedRight\arraybackslash\ttfamily\scriptsize}p{0.935\linewidth}@{}}
1 & \# 周芷若 · 终章 · 白莲重生\\
2 & \mbox{}\\
3 & > **Category**: 角色肖像 · 周芷若\\
4 & > **Level**: 18/18\\
5 & > **Frame**: 角色肖像\\
6 & > **Aspect Ratio**: 16:9\\
7 & > **Art Style**: Chinese Gongbi (工笔) figure painting — fine brushwork, detailed outlines, layered colors, precise facial features and clothing patterns for lotus (the most refined flower painting subject), Chinese Xieyi (写意) ink painting — freehand brushwork, expressive washes, atmospheric, poetic suggestion over literal depiction for distant mountains, fusion of flower-bird painting and landscape, spiritual fulfillment\\
8 & \mbox{}\\
9 & ---\\
10 & \mbox{}\\
11 & \#\# \promptemoji{📖} 场景描述 (Scene Description)\\
12 & \mbox{}\\
13 & 峨眉山下一方莲池。盛夏。满池白莲盛开——有的完全绽放如白玉盘，有的含苞待放如毛笔头，有的已结成莲蓬（莲子饱满）。碧绿的荷叶如伞盖层层叠叠——叶面上滚动着晶莹的水珠（清晨的露水）。池水清澈见底——可见几尾红鲤在莲茎之间穿游。周芷若（约三十五岁）赤足站在池边浅水中——身着素白麻布衣（最简单最朴素的衣料——麻的质地可见），长发随意绾在脑后。她弯下腰——双手轻轻捧起池水洗面。水从她指缝流下——在晨光中如碎银。她的脸上有水珠——分不清是池水还是别的什么。但她在笑——嘴角微微上扬，眼睛弯成月牙。她身后是峨眉山的晨曦——金顶在远处闪光。莲池边有一座小茅亭——亭中有石桌石凳，桌上放着笔墨和一册书稿。白莲花在她周围——仿佛她也是其中一朵。\\
14 & \mbox{}\\
15 & ---\\
16 & \mbox{}\\
17 & \#\# \promptemoji{🖼} 构图指南 (Composition Guide)\\
18 & \mbox{}\\
19 & 环绕式构图。周芷若在画面中央偏下，周围的莲叶和莲花以她为圆心呈弧形分布——形成自然的"花环"构图。她的弯腰动作形成柔和的曲线。池水波纹从她脚下向外扩散（同心圆）。峨眉山在远景上方——天地轴线将她、莲池、金顶连成一线。小茅亭在画面左侧——平衡构图。\\
20 & \mbox{}\\
21 & ---\\
22 & \mbox{}\\
23 & \#\# \promptemoji{💡} 光影设计 (Lighting Design)\\
24 & \mbox{}\\
25 & soft rose-gold morning light, long shadows, fresh and hopeful, mist rising。清晨的阳光从峨眉山方向（画面右上）斜射入莲池——在每一片莲叶上形成金色的边缘光。水珠在阳光下闪烁如钻石。周芷若面部被柔和晨光照亮——水光的反射在她脸上形成流动的光影（动态光影，来自水面的波光）。\\
26 & \mbox{}\\
27 & ---\\
28 & \mbox{}\\
29 & \#\# \promptemoji{🎨} 色彩方案 (Color Palette)\\
30 & \mbox{}\\
31 & 主调：莲叶翠绿 + 白莲花乳白（花心淡黄）+ 池水清澈淡蓝绿。人物：素白麻衣（偏暖白——麻的本色）+ 肤色健康暖象牙。红鲤：朱砂红（水中一点亮色）。峨眉山：远景青蓝。晨光：暖金色。整体清新、纯净、充满生机——这是"重生"的颜色。不再是"白衣"的冷酷或"黑衣"的阴郁——是"麻衣"的质朴与真实。\\
32 & \mbox{}\\
33 & ---\\
34 & \mbox{}\\
35 & \#\# \promptemoji{🔑} 关键元素 (Key Elements)\\
36 & \mbox{}\\
37 & 白莲（不同阶段——花苞、初放、盛放、莲蓬——生命的完整循环），莲叶（碧绿如盖，叶脉清晰，水珠滚动），池水（清澈见底，红鲤游弋），赤足（站在浅水中，脚踝以下没入水中），掬水（双手捧水，水从指缝成串落下——水珠在晨光中清晰可见），麻衣（粗麻布纹理可见——最质朴的衣料），微笑（嘴角自然上扬，不是刻意——是内心的平静自然流露），峨眉金顶（远景——在她生命中既是牢笼也是归属的地方，如今只是风景），茅亭（简朴木构，桌上有书稿——《峨眉女科》）。\\
38 & \mbox{}\\
39 & ---\\
40 & \mbox{}\\
41 & \#\# \promptemoji{😌} 情绪氛围 (Mood \& Atmosphere)\\
42 & \mbox{}\\
43 & 重生。她走过黑暗，走过迷失，走过痛苦——最终成了一朵白莲。出淤泥而不染，濯清涟而不妖——不只是花，也是她。\\
44 & \mbox{}\\
45 & \mbox{}\\
46 & \#\# \promptemoji{🎨} 通用风格指南 (Universal Style Guide)\\
47 & \mbox{}\\
48 & \#\#\# 核心美术风格 (Core Art Style)\\
49 & - **Primary**: 金庸武侠水墨风 (Jin Yong Wuxia Ink Painting Aesthetic)\\
50 & - **Figure Style**: 中国工笔人物 (Chinese Gongbi Figure Painting) — fine, precise brushwork for faces, hands, and clothing details\\
51 & - **Background Style**: 中国写意山水 (Chinese Xieyi Landscape) — expressive, atmospheric ink washes for environments\\
52 & - **Fusion Approach**: Gongbi precision in foreground subjects + Xieyi freedom in backgrounds = balanced classical Chinese aesthetic\\
53 & \mbox{}\\
54 & \#\#\# 时代考据 (Historical Accuracy)\\
55 & - **Period**: 元末明初 (Late Yuan / Early Ming Dynasty, c. 1350-1400)\\
56 & - **Architecture**: Yuan Dynasty official buildings (Dougong brackets, hip-and-gable roofs, colored glazed tiles), Mongolian yurts (gers), Song-style garden pavilions\\
57 & - **Clothing**:\\
58 &   - Mongolian aristocracy: White sable furs, gold ornaments, leather boots, rounded-collar robes (质孙服)\\
59 &   - Han Chinese women: Cross-collar ruqun (襦裙), beizi (褙子) outer jackets, flowing sleeves\\
60 &   - Martial artists: Tight-cuffed training clothes, cloth boots, minimal accessories\\
61 & - **Weapons**: Historically accurate Chinese swords (jian 剑), sabers (dao 刀), spears (qiang 枪)\\
62 & \mbox{}\\
63 & \#\#\# 构图原则 (Composition Principles)\\
64 & - **Rule of Thirds**: Primary subject at 1/3 or 2/3 position, not dead center\\
65 & - **Negative Space** (留白): Intentional empty areas for visual breathing room — a core Chinese painting principle\\
66 & - **Depth Layers**: Foreground (details/framing) → Middle ground (main action) → Background (atmosphere)\\
67 & - **Leading Lines**: Natural elements (rivers, paths, architectural lines) guiding eye to focal point\\
68 & - **Scale Reference**: Include recognizable elements (trees, buildings, figures) for size context\\
69 & \mbox{}\\
70 & \#\#\# 色彩体系 (Color System)\\
71 & - **Warm Earth Base**: 赭石 (burnt sienna), 土黄 (yellow ochre), 熟褐 (burnt umber)\\
72 & - **Jade Accents**: 石青 (azurite blue), 石绿 (malachite green) for landscapes\\
73 & - **Dramatic Reds**: 朱砂 (cinnabar red), 胭脂 (rouge) for important moments\\
74 & - **Ink Blacks**: 焦墨 (burnt ink), 浓墨 (dark ink), 淡墨 (light ink) — five gradations of black\\
75 & - **Gold Highlights**: 泥金 (gold paint) for divine light, sacred objects, sun rays\\
76 & - **Avoid**: Neon colors, modern synthetic palettes, pure RGB primaries\\
77 & \mbox{}\\
78 & \#\#\# 光影处理 (Lighting Treatment)\\
79 & - Classical Chinese painting traditionally uses flat/diffuse lighting\\
80 & - For this project: Blend traditional flat lighting with subtle directional light\\
81 & - Strong directional light ONLY in dramatic scenes (fire, divine intervention)\\
82 & - Prefer: Soft, diffuse illumination that reveals form without harsh shadows\\
83 & - Light source: Usually from upper-left or upper-right, gentle falloff\\
84 & \mbox{}\\
85 & \#\#\# 技术参数 (Technical Specifications)\\
86 & - **Aspect Ratio**: 16:9 (landscape orientation)\\
87 & - **Resolution Target**: 1792×1024 pixels (DALL-E 3) or equivalent\\
88 & - **Image Type**: Digital painting in classical Chinese style\\
89 & - **Level of Detail**: Museum-quality fine art reproduction level\\
90 & \mbox{}\\
91 & \#\#\# 绝对禁止 (Strictly Forbidden)\\
92 & - \promptemoji{❌} Modern clothing, accessories, or technology (no smartphones, modern glasses, etc.)\\
93 & - \promptemoji{❌} Western fantasy aesthetics (no European castles, dragons, knights)\\
94 & - \promptemoji{❌} Anime/manga style (no large eyes, simplified cell shading)\\
95 & - \promptemoji{❌} Photorealism (this is a PAINTING, not a photograph)\\
96 & - \promptemoji{❌} Excessive gore/violence (suggest rather than show graphic content)\\
97 & - \promptemoji{❌} Text/watermarks/signatures (clean image, no overlaid text)\\
98 & - \promptemoji{❌} 3D rendered appearance\\
99 & \mbox{}\\
100 & \mbox{}\\
101 & ---\\
102 & \mbox{}\\
103 & \#\# \promptemoji{🤖} GPT生成提示词 (GPT Generation Prompt)\\
104 & \mbox{}\\
105 & \#\#\# 中文提示词 (Chinese Prompt)\\
106 & ```\\
107 & 周芷若 · 终章 · 白莲重生。峨眉山下一方莲池。盛夏。满池白莲盛开——有的完全绽放如白玉盘，有的含苞待放如毛笔头，有的已结成莲蓬（莲子饱满）。碧绿的荷叶如伞盖层层叠叠——叶面上滚动着晶莹的水珠（清晨的露水）。池水清澈见底——可见几尾红鲤在莲茎之间穿游。周芷若（约三十五岁）赤足站在池边浅水中——身着素白麻布衣（最简单最朴素的衣料——麻的质地可见），长发随意绾在脑后。她弯下腰——双手轻轻捧起池水洗面。水从她指缝流下——在晨光中如碎银。她的脸上有水珠——分不清是池水还是别的什么。但她在笑——嘴角微微上扬，眼睛弯成月牙。她身后是峨眉山的晨曦——金顶在远处闪光。莲池边有一座小茅亭——亭中有石桌石凳，桌上放着笔墨和一册书稿。白莲花在她周围——仿佛她也是其中一朵。 构图：环绕式构图。周芷若在画面中央偏下，周围的莲叶和莲花以她为圆心呈弧形分布——形成自然的"花环"构图。她的弯腰动作形成柔和的曲线。池水波纹从她脚下向外扩散（同心圆）。峨眉山在远景上方——天地轴线将她、莲池、金顶连成一线。小茅亭在画面左侧——平衡构图。 光影：soft rose-gold morning light, long shadows, fresh and hopeful, mist rising。清晨的阳光从峨眉山方向（画面右上）斜射入莲池——在每一片莲叶上形成金色的边缘光。水珠在阳光下闪烁如钻石。周芷若面部被柔和晨光照亮——水光的反射在她脸上形成流动的光影（动态光影，来自水面的波光）。 色彩：主调：莲叶翠绿 + 白莲花乳白（花心淡黄）+ 池水清澈淡蓝绿。人物：素白麻衣（偏暖白——麻的本色）+ 肤色健康暖象牙。红鲤：朱砂红（水中一点亮色）。峨眉山：远景青蓝。晨光：暖金色。整体清新、纯净、充满生机——这是"重生"的颜色。不再是"白衣"的冷酷或"黑衣"的阴郁——是"麻衣"的质朴与真实。 风格：Chinese Gongbi (工笔) figure painting — fine brushwork, detailed outlines, layered colors, precise facial features and clothing patterns for lotus (the most refined flower painting subject), Chinese Xieyi (写意) ink painting — freehand brushwork, expressive washes, atmospheric, poetic suggestion over literal depiction for distant mountains, fusion of flower-bird painting and landscape, spiritual fulfillment，中国工笔人物画结合写意山水背景，金庸武侠美学，元代历史考据，博物馆级精细度，16:9横构图。\\
108 & ```\\
109 & \mbox{}\\
110 & \#\#\# English Prompt\\
111 & ```\\
112 & Zhou Zhiruo character portrait 18: 周芷若 · 终章 · 白莲重生. 峨眉山下一方莲池。盛夏。满池白莲盛开——有的完全绽放如白玉盘，有的含苞待放如毛笔头，有的已结成莲蓬（莲子饱满）。碧绿的荷叶如伞盖层层叠叠——叶面上滚动着晶莹的水珠（清晨的露水）。池水清澈见底——可见几尾红鲤在莲茎之间穿游。周芷若（约三十五岁）赤足站在池边浅水中——身着素白麻布衣（最简单最朴素的衣料——麻的质地可见），长发随意绾在脑后。她弯下腰——双手轻轻捧起池水洗面。水从她指缝流下——在晨光中如碎. 峨眉山下一方莲池。盛夏。满池白莲盛开——有的完全绽放如白玉盘，有的含苞待放如毛笔头，有的已结成莲蓬（莲子饱满）。碧绿的荷叶如伞盖层层叠叠——叶面上滚动着晶莹的水珠（清晨的露水）。池水清澈见底——可见几尾红鲤在莲茎之间穿游。周芷若（约三十五岁）赤足站在池边浅水中——身着素白麻布衣（最简单最朴素的衣料——麻的质地可见），长发随意绾在脑后。她弯下腰——双手轻轻捧起池水洗面。水从她指缝流下——在晨光中如碎银。她的脸上有水珠——分不清是池水还是别的什么。但她在笑——嘴角微微上扬，眼睛弯成月牙。她身后是峨眉山的晨曦——金顶在远处闪光。莲池边有一座小茅亭——亭中有石桌石凳，桌上放着笔墨和一册书稿。白莲花在她周围——仿佛她也是其中一朵。 Composition: 环绕式构图。周芷若在画面中央偏下，周围的莲叶和莲花以她为圆心呈弧形分布——形成自然的"花环"构图。她的弯腰动作形成柔和的曲线。池水波纹从她脚下向外扩散（同心圆）。峨眉山在远景上方——天地轴线将她、莲池、金顶连成一线。小茅亭在画面左侧——平衡构图。 Lighting: soft rose-gold morning light, long shadows, fresh and hopeful, mist rising。清晨的阳光从峨眉山方向（画面右上）斜射入莲池——在每一片莲叶上形成金色的边缘光。水珠在阳光下闪烁如钻石。周芷若面部被柔和晨光照亮——水光的反射在她脸上形成流动的光影（动态光影，来自水面的波光）。 Color palette: 主调：莲叶翠绿 + 白莲花乳白（花心淡黄）+ 池水清澈淡蓝绿。人物：素白麻衣（偏暖白——麻的本色）+ 肤色健康暖象牙。红鲤：朱砂红（水中一点亮色）。峨眉山：远景青蓝。晨光：暖金色。整体清新、纯净、充满生机——这是"重生"的颜色。不再是"白衣"的冷酷或"黑衣"的阴郁——是"麻衣"的质朴与真实。 Style: Chinese Gongbi (工笔) figure painting — fine brushwork, detailed outlines, layered colors, precise facial features and clothing patterns for lotus (the most refined flower painting subject), Chinese Xieyi (写意) ink painting — freehand brushwork, expressive washes, atmospheric, poetic suggestion over literal depiction for distant mountains, fusion of flower-bird painting and landscape, spiritual fulfillment, Chinese gongbi figure painting with xieyi landscape background, Jin Yong wuxia aesthetic, Yuan Dynasty historical accuracy, museum-quality fine art, 16:9 horizontal composition, highly detailed, classical Chinese painting technique.\\
113 & ```\\
114 & \mbox{}\\
115 & ---\\
116 & \mbox{}\\
117 & \#\# \promptemoji{🚫} 负面提示词 (Negative Prompt)\\
118 & ```\\
119 & modern clothing, modern technology, smartphones, cars, airplanes, western architecture, european castles, anime style, manga style, cartoon, cel shading, large eyes, photorealistic, 3D render, photograph, text overlay, watermark, signature, excessive gore, graphic violence, neon colors, synthetic palette, abstract, minimalist, blurry, low resolution, distorted faces, extra limbs, bad anatomy\\
120 & \mbox{}\\
121 & ```\\
122 & \mbox{}\\
123 & ---\\
124 & \mbox{}\\
125 & \#\# \promptemoji{📐} 技术参数 (Technical Specs)\\
126 & - **Generator**: GPT-4o with Image Generation / DALL·E 3\\
127 & - **Aspect Ratio**: 16:9\\
128 & - **Style**: vivid (DALL·E 3) / natural\\
129 & - **Quality**: HD\\
130 & - **Expected Resolution**: 1792×1024 or higher\\
131 & \mbox{}\\
132 & ---\\
133 & \mbox{}\\
134 & *Generated for 浙江大学《中西医学与美学》期末大作业 · 倚天屠龙记·胡笳十八拍*\\
\end{longtable}
\endgroup
